% GENERATED FILE. Edit canonical lessons, then run npm run generate:books.
% canonical-source-hash: fnv1a64:6cb8f4e02c6e1e7f
% canonical-lessons: GU-W11-a2-connected-composition

\chapter{A connected A2 message}
\label{ch:gu-a-connected-a2-message}
\hlchaptermodality{144}

\begin{chapteropening}
\textbf{By the end of this chapter:} \emph{I can plan and write a connected 70–90 word Gujarati message in Gujarati script for a named reader and practical purpose.}

Complete chapter 144 by writing a model-free 70–90 word connected Gujarati message for a named reader and practical purpose.
\end{chapteropening}

\section[(connected A2 writing)]{A connected message about familiar life}

\label{lesson:GU-W11-a2-connected-composition}



\begin{quote}
A connected message is not a pile of vocabulary and not one long sentence. It is a short route for one reader. Before writing, draw eight empty boxes and give them these jobs:

\begin{enumerate}
\raggedright
  \item reader and greeting;
  \item reason for writing;
  \item your name and where you live;
  \item a language you speak and work or activity you do;
  \item one more familiar detail;
  \item one preference and one reason;
  \item one direct question for the reader;
  \item a close.
\end{enumerate}

Write only one or two English cue words in each box. Do not draft Gujarati yet.
\end{quote}

\subsection*{Writing — connected composition}
There is no model, sentence bank, or copyable answer on this page.

\textbf{Reader:} your friend Kiran. \textbf{Purpose:} help Kiran understand your familiar life and choose an activity you could do together.

Write \textbf{70–90 words in Gujarati using Gujarati script}. Use every planning box. Include your name, where you live, a language, at least two familiar actions or activities, one preference, one reason, and one direct question for Kiran. Make every detail help with the shared activity. The instructions give content, not sentences: choose only words and forms this book has taught, and connect them yourself. If you cannot form a line honestly, make it simpler; do not borrow a Hindi form to fill the gap.

This first connected attempt is untimed. Put away romanization, dictionaries, translators, generative assistants, and spell-checkers. Do not copy from an earlier lesson. Ordinary legible Gujarati handwriting is enough.

When the whole message exists, count its words once. If it is short, add one useful taught detail to a planning box. If it is long, remove a detail that does not help Kiran. Do not replace the message with a memorised model.

\begin{tcolorbox}[breakable,colback=teal!4,colframe=teal!35!black,title={Before you move on}]
Review the same attempt in four separate passes:

\begin{enumerate}
\raggedright
  \item \textbf{Task fulfilment and relevant content} — Kiran, the familiar-life purpose, shared activity, and every requested content point are present.
  \item \textbf{Comprehensibility and organisation} — the eight jobs form a route rather than a list of unrelated words or lines.
  \item \textbf{Gujarati vocabulary and grammar} — every form comes from the taught Gujarati track; a familiar shorter line is better than an invented form.
  \item \textbf{Orthographic control} — Gujarati letters, vowel signs, word boundaries, and punctuation remain readable.
\end{enumerate}

Mark one strength and one next step. Repair only the next step; keep the original attempt so the next 70–90 word message can show real growth.
\end{tcolorbox}
